%-------------------------------------------------------------------
%  Рядовой Тимофей — резюме (русская версия)
%  Компиляция:  pdflatex resume-ru.tex   (нужен пакет cm-super для кириллицы)
%  Затем положи resume-ru.pdf в ../public/
%-------------------------------------------------------------------
\documentclass[a4paper,11pt]{article}

\usepackage{latexsym}
\usepackage[empty]{fullpage}
\usepackage{titlesec}
\usepackage{marvosym}
\usepackage[usenames,dvipsnames]{color}
\usepackage{verbatim}
\usepackage{enumitem}
\usepackage[pdftex]{hyperref}
\usepackage{fancyhdr}

% --- поддержка кириллицы (pdflatex) ---
\usepackage{cmap}
\usepackage[T2A]{fontenc}
\usepackage[utf8]{inputenc}
\usepackage[english,russian]{babel}

\pagestyle{fancy}
\fancyhf{} % clear all header and footer fields
\fancyfoot{}
\renewcommand{\headrulewidth}{0pt}
\renewcommand{\footrulewidth}{0pt}

% Adjust margins
\addtolength{\oddsidemargin}{-0.375in}
\addtolength{\evensidemargin}{-0.375in}
\addtolength{\textwidth}{1in}
\addtolength{\topmargin}{-.5in}
\addtolength{\textheight}{1.0in}

\setlength{\footskip}{5pt}

\urlstyle{same}

\raggedbottom
\raggedright
\setlength{\tabcolsep}{0in}

% Sections formatting
\titleformat{\section}{
  \vspace{-4pt}\scshape\raggedright\large
}{}{0em}{}[\color{black}\titlerule \vspace{-5pt}]

%-------------------------
% Custom commands
\newcommand{\resumeItem}[2]{
  \item\small{
    \textbf{#1}{: #2 \vspace{-2pt}}
  }
}

\newcommand{\resumeSubheading}[4]{
  \vspace{-1pt}\item
    \begin{tabular*}{0.97\textwidth}{l@{\extracolsep{\fill}}r}
      \textbf{#1} & #2 \\
      \textit{\small#3} & \textit{\small #4} \\
    \end{tabular*}\vspace{-5pt}
}

\newcommand{\resumeSubItem}[2]{\resumeItem{#1}{#2}\vspace{-4pt}}

\renewcommand{\labelitemii}{$\circ$}

\newcommand{\resumeSubHeadingListStart}{\begin{itemize}[leftmargin=*]}
\newcommand{\resumeSubHeadingListEnd}{\end{itemize}}
\newcommand{\resumeItemListStart}{\begin{itemize}}
\newcommand{\resumeItemListEnd}{\end{itemize}\vspace{-5pt}}

%-------------------------------------------
%%%%%%  РЕЗЮМЕ  %%%%%%%%%%%%%%%%%%%%%%%%%%%%

\begin{document}

%----------ШАПКА-----------------
\begin{tabular*}{\textwidth}{l@{\extracolsep{\fill}}r}
  \textbf{\Large Рядовой Тимофей Сергеевич} & Почта : \href{mailto:timaryadovou@yandex.ru}{timaryadovou@yandex.ru}\\
  \href{https://timryadovouu.github.io}{timryadovouu.github.io} & Телефон : +7 999 035-03-29 \\
  \href{https://t.me/urokodaki_sakonji}{telegram} & \\
\end{tabular*}


%-----------ОБРАЗОВАНИЕ-----------------
\section{Образование}
  \resumeSubHeadingListStart
    \resumeSubheading
      {СПбГЭТУ «ЛЭТИ»}{Санкт-Петербург}
      {Магистратура — «Системы и технологии технического зрения»}{Сен. 2026 -- н.в.}
    \resumeSubheading
      {Университет ИТМО}{Санкт-Петербург}
      {Бакалавриат — «Конструирование и технология электронных средств»}{Сен. 2021 -- Июнь 2026}
  \resumeSubHeadingListEnd


%-----------ОПЫТ РАБОТЫ-----------------
\section{Опыт работы}
  \resumeSubHeadingListStart

    \resumeSubheading
      {ПАО «Совкомбанк»}{Санкт-Петербург}
      {Аналитик данных}{Сен. 2025 -- н.в.}
      \resumeItemListStart
        \resumeItem{Пайплайны}
          {Формирование конвейеров для автоматизации обработки и трансформации данных перед анализом.}
        \resumeItem{Выгрузка данных}
          {Выгрузка данных из баз данных, хранилищ, витрин и других источников.}
        \resumeItem{Подготовка данных}
          {Предобработка, очистка и подготовка данных для дальнейшего анализа.}
        \resumeItem{Презентации}
          {Подготовка и проведение презентаций с результатами анализа и моделирования.}
        \resumeItem{Отчётность}
          {Построение и поддержка регулярной отчётности и дашбордов (Power BI).}
        \resumeItem{Риск-тестирование}
          {Риск-тестирование кредитных портфелей банка методом Монте-Карло.}
      \resumeItemListEnd

  \resumeSubHeadingListEnd


%-----------КУРСЫ-----------------
\section{Курсы и тренинги}
  \resumeSubHeadingListStart
    \resumeSubheading
      {YADRO AI School}{Университет ИТМО}
      {9-месячная программа по ИИ и машинному обучению}{2024 -- 2025}
    \resumeSubheading
      {Основы сетевых технологий}{Университет ИТМО и D-Link}
      {4-месячная программа по компьютерным сетям}{2024}
  \resumeSubHeadingListEnd


%-----------ПРОЕКТЫ-----------------
\section{Проекты}
  \resumeSubHeadingListStart
    \resumeSubItem{\href{https://github.com/timryadovouu/mac-notch}{mac-notch}}
      {Нативная утилита для macOS, превращающая вырез камеры MacBook в интерактивный хаб в духе Dynamic Island: таймер, буфер обмена, управление музыкой, to-do и учёт экранного времени. Swift, SwiftUI, AppKit.}
    \resumeSubItem{Почему сталкивающиеся блоки образуют число Пи}
      {Выступление на студенческой весне ИТМО о механике блоков, число столкновений которых даёт цифры числа Пи.}
    \resumeSubItem{Кудиты в квантовых пространствах высокой размерности}
      {Выступление на квантовом семинаре ИТМО о кудитах и квантовых системах высокой размерности.}
  \resumeSubHeadingListEnd


%--------НАВЫКИ------------
\section{Навыки}
 \resumeSubHeadingListStart
   \item{
     \textbf{Языки программирования}{: Python, C, JavaScript / TypeScript, SQL, Bash}
   }
   \item{
     \textbf{Библиотеки и фреймворки}{: NumPy, Pandas, Scikit-learn, PyTorch, FastAPI, Django, React}
   }
   \item{
     \textbf{Инструменты}{: Git, Docker, Jupyter Notebook, Linux}
   }
 \resumeSubHeadingListEnd


%--------ДОПОЛНИТЕЛЬНО------------
\section{Дополнительно}
 \resumeSubHeadingListStart
   \item{
     \textbf{Языки}{: русский (родной), английский (B2)}
   }
   \item{
     \textbf{Интересы}{: алгоритмы, математика, физика, машинное обучение, NLP, научная фантастика; фехтование, бадминтон, велоспорт, настольный теннис, фортепиано и гитара}
   }
 \resumeSubHeadingListEnd


%-------------------------------------------
\end{document}
